\chapter{Syntax}
\label{ch:syntax}

This section mainly focuses on the syntax rules of the grammar; for more
details about a specific feature, refer to the dedicated section.

\vspace{\baselineskip}

A valid Giovi program starts with zero or more top-level statements; there are
11 valid statements:

\begin{minted}{antlr}
program: statement* ;
statement
	: statement_declaration
	| statement_return
	| statement_invocation
	| statement_redirection
	| statement_import
	| statement_if
	| statement_break
	| statement_skip
	| statement_repeat
	| statement_for
	;

statement_declaration
	: statement_declaration_alias
	| statement_declaration_module
	| statement_declaration_struct
	| statement_declaration_constant
	| statement_declaration_variable
	| statement_declaration_function
	;
\end{minted}

In many rules, either a single statement or a block of statements are allowed;
the following "fragment" rules will be used.

\begin{minted}{antlr}
block: '{' statement* '}';
block_or_statement: block | statement ;
\end{minted}

\paragraph{Identifiers} The \code{identifier} rule describes the allowed
identifiers to be used as names for variables, constants, types etc.

\begin{minted}{antlr}
identifier: [a-zA-Z] [a-zA-Z0-9_]* ;
\end{minted}

It must start with a letter and can contain any letter, number or underscore.

\clearpage

\section{Alias declaration}

\begin{minted}{antlr}
statement_declaration_alias: 'alias' expression 'as' identifier ;
\end{minted}

Sometimes a type, a function, a constant or something else may be named in a
too verbose way. You can use the alias statement to basically refer to the same
thing but with a different name.

\begin{minted}{\giovilexer}
alias AVeryVerboseAndLongStructName as SomeStruct
\end{minted}

\section{Import statement}

\begin{minted}{antlr}
statement_import
	: 'import' import_entry
	| 'import' '[' import_entry+ ']'
	;
import_entry: import_uri ('.' import_uri)* ('as' alias=identifier)? ;
import_uri
	: identifier
	| literal_normal_text
	| literal_exact_text
	;
\end{minted}

The import statement is used to import one or more modules
(\see{ch:modules}).

\begin{minted}{\giovilexer}
import giovi.util.Random
giovi.util.Random.integer() -> _out
\end{minted}

To avoid fully qualifying the module every time, it can be aliased. It's a
shortcut for importing and then aliasing with two different statements.

\begin{minted}{\giovilexer}
import giovi.util.Random as Random
Random.integer() -> _out
\end{minted}

For more details, see~\See{sec:import}.

\section{Module declaration}

\begin{minted}{antlr}
statement_declaration_module
	: 'module' identifier body=block
	;
\end{minted}

A module\footnotemark[1] is a "container" for code, and shares the same behavior of a scope: it
can hold code, variables, functions, constants, types definition etcetera; in
fact it can be seen as a "named scope". See \See{ch:modules}.

\footnotetext[1]{
	Modules are still considered an experimental feature and, as such, the
	syntax is considered unstable.
}

\section{Constant declaration}

\begin{minted}{antlr}
statement_declaration_constant
	: 'const' identifier type? '<-' expression
	;
\end{minted}

Constants are read-only values; their value is defined at declaration and
cannot change.\\*
The type can be omitted as it will be inferred from the expression.\\*
By convention, constants identifiers are written in \code{SCREAMING\_CASE}.

\todo{decidere se usare shallow o deep immutability}

\begin{minted}{\giovilexer}
const PI Real <- 3.14
const TAU <- PI * 2
\end{minted}

\section{Variable declaration}

\begin{minted}{antlr}
statement_declaration_variable
	: 'var' identifier type? ('default' expression)? '<-' expression
	| identifier type ('default' expression)?
	;
\end{minted}

A variable is identified by a name, and it is of a specific type.\\*
By convention, variable identifiers are written in \code{pascalCase}.\\*
Using the \code{var} keyword allows you to omit the type and let it be inferred
from the assigned expression.

\begin{minted}{\giovilexer}
myName Text <- "Simone"
var myNickname <- "simonvic"
\end{minted}

\section{Return statement}

\begin{minted}{antlr}
statement_return
	: ('return' | '=>') expression?
	;
\end{minted}

The return statement can only be used within a function scope.

\clearpage

\section{Function declaration}

\begin{minted}{antlr}
statement_declaration_function
	: identifier '(' formal_parameters? ')' ('=>' return_type=type)?
	  body=block
	;

formal_parameters
	: formal_parameter (',' formal_parameter)*
	;

formal_parameter
	: (label=identifier ':')? name=identifier type ('default' expression)?
	;
\end{minted}

Functions are code blocks that can be invoked at a later time.\\*
By convention, function identifiers are written in \code{pascalCase}.

\begin{minted}{\giovilexer}
printHelloWorld() {
	"Hello World" -> _out
}

greetPerson(name Text, surname Text) {
	"Greetings $name, $surname" -> _out
}

getFunLanguage() => Text {
	=> "Giovi"
}
\end{minted}

For more details, see~\See{ch:functions}.

\section{Invocation statement}

\begin{minted}{antlr}
statement_invocation
	: expression '(' actual_parameters ')'
	;

actual_parameters
	: actual_parameter (',' actual_parameter)*
	;

actual_parameter
	: (label=identifier ':')? expression
	;
\end{minted}

When invoking a function, arguments can be passed using a traditional
positional approach or a nominal approach\footnotemark[1].

\footnotetext[1]{
	Albeit the grammar doesn't enforce one approach or the other, the current
	implementation supports either all arguments to be passed as positional or
	as nominal. For more details, see~\See{ch:functions}.
}

\section{Redirection statement}

\begin{minted}{antlr}
statement_redirection
	: statement_redirection_rtl
	| statement_redirection_ltr
	;

statement_redirection_rtl
	: dest=pattern_redir '<-' src=expression
	| dest=expression ('<-', '+<-', '-<-', '*<-', '/<-') src=expression
	;

statement_redirection_ltr
	: src=expression '->' dest=pattern_redir
	| src=expression ('->', '->+', '->-', '->*', '->/') content=expression
	;

pattern_redir
	: pattern_placeholder
	| pattern_unary
	| pattern_placeholder
	;
\end{minted}

The redirection statement is simply an assignment of a value to
\textit{something}.

A peculiarity of the statement is that it can be written "left to right" or
"right to left". There is no difference between the two directions, except for
readability. Sometimes it's just better to use one instead of the other.

\todo{link a FAQ}

\begin{minted}{\giovilexer}
"hello world" -> _out     # Printing
counter +<- 1             # Increasing counter by one
acceleration ->+ velocity # Increasing velocity by "acceleration"
\end{minted}

It's common having to perform operations such as sums or concatenations with a
variable, and then to apply the resulting value of the operation to the same
variable. Giovi gives you a shorthand for such scenarios: it's possible to
append or suffix the \code{+}, \code{-}, \code{*} and \code{/} operators to the
redirection operator (both left-to-right and right-to-left variant); when doing
so, the specified operation will be performed using the destination of the
assignment as the left-hand-side operand and the expression as the
right-hand-side operand; finally, the result will be applied to the
destination.

Beware that the result will be identical if you were to do the operation in a
separate expression and then assigning the result, meaning that if the operator
is not supported or the result of the expression is incompatible, the
assignment will fail.

\begin{minted}{\giovilexer}
counter Real <- 6.9
counter *<- 10 # multiply by 10

funnyNumbers [Int]
funnyNumbers +<- 69 # insert 69
funnyNumbers +<- [420, 67] # concatenate

sentence Text <- "hello world how are you?"
sentence /<- " " # ERROR: can't assign [hello, world, how, are, you?]
\end{minted}

\clearpage

\section{If statement}

\begin{minted}{antlr}
statement_if
	: statement_simple_if
	| statement_multi_if
	;
\end{minted}

The \code{if} statement is used to evaluate boolean expressions exclusively.\\*
There is a "traditional" \code{if}, where a single expression is evaluated and
the appropriate execution flow is chosen, and a "multi-if" version.

\subsection{Simple If statement}

\begin{minted}{antlr}
statement_simple_if
	: 'if' expression 'do' statement_or_block ('else' statement_or_block)?
	;
\end{minted}

\begin{minted}{\giovilexer}
if someConditionIsTrue() do something() else somethingElse()
if someConditionIsTrue() do {
	something()
} else {
	somethingElse()
}
\end{minted}

\subsection{Multi If statement}

\begin{minted}{antlr}
statement_multi_if
	: 'if' subject=expression multi_if_body ('else' statement_or_block)?
	| 'if' subjectless_multi_if_body ('else' statement_or_block)?
	;
\end{minted}

Very often it is needed to check multiple conditions sequentially.

\begin{minted}{\giovilexer}
if name = "Simone" do {
	"hey it's Simone!" -> _out
} else if name = "Mario" do {
	"heyo Mario" -> _out
	"where is Luigi?" -> _out
} else if name = "Luigi" do {
	"heyo green Mario!" -> _out
}
\end{minted}

To improve readability, the \code{if} statement allows you to write the same in
a more concise way.

\begin{minted}{\giovilexer}
if {
	name = "Simone" do "hey it's Simone!" -> _out
	name = "Mario" do {
		"heyo Mario" -> _out
		"where is Luigi?" -> _out
	}
	name = "Luigi" do "heyo green Mario!" -> _out
}
\end{minted}

\nb{Any boolean expression is allowed and exhaustiveness is not enforced! Be
mindful of the branch order; since they're evaluated in the same order they're
written, some branch may become unreachable if it's dominated by a previous
one.}

\paragraph{Subject} More often than not, the conditions involve the same
"subject". The \code{if} statement allows you to write it just once and then
omit it in the left-hand side of the condition

\begin{minted}{\giovilexer}
if name {
	= "Simone" do "hey it's Simone!" -> _out
	= "Mario" do {
		"heyo Mario" -> _out
		"where is Luigi?" -> _out
	}
	= "Luigi" do "heyo green Mario!" -> _out
}
\end{minted}

Common boolean operators are allowed

\begin{minted}{\giovilexer}
if getSomeNumber() {
	= 69 or = 420 do "Nice!" -> _out
	> 100 and < 999 do "big number" -> _out
}
\end{minted}

\paragraph{Branch guard} A branch can also be "guarded" with another (possibly
unrelated) boolean condition

\begin{minted}{\giovilexer}
var funnyNumbers <- [69, 420] 
if getSomeNumber() {
	# if equal to 69, but only if someCondition() is true
	= 69 when someCondition() do "Nice!" -> _out

	# if is contained in funnyNumbers list,
	# but only when 67 is not contained in funnyNumbers
	in funnyNumbers when 67 not in funnyNumbers
}
\end{minted}

\clearpage

\section{Break and Skip statement}

Both statements are used to control the flow of a loop, therefore can only be
used within a loop.

\subsection{Break statement}

\begin{minted}{antlr}
statement_break: 'break';
\end{minted}

The break statement will stop the loop and the execution will continue with the
next instruction.

\subsection{Skip statement}

\begin{minted}{antlr}
statement_skip: 'skip';
\end{minted}

The skip statement will stop the current loop iteration and skip to the next
iteration.

\section{Repeat statement}

\begin{minted}{antlr}
statement_repeat: 'repeat' ('if' expression)? statement_or_block ;
\end{minted}

The repeat statement allows you to repeat a statement or block of code. It
allows an optional condition that gets checked at the beginning of every loop;
if the condition fails, the iteration stops.

\begin{minted}{\giovilexer}
# this code will run indefinitely
repeat {
	"hello world" -> _out
}
\end{minted}

\begin{minted}{\giovilexer}
# the loop continues while someCondition is true
repeat if someCondition() {
	"hello world" -> _out
}
\end{minted}

A "do-while" like statement, can be achieved with a \code{break} statement.

\begin{minted}{\giovilexer}
# the loop continues until someCondition is true
repeat {
	"hello world" -> _out
	if someCondition() do break
}
\end{minted}

\clearpage

\section{For statement}

\begin{minted}{antlr}
statement_for
	: 'for' pattern_iterator 'in' iterable 'do' statement_or_block
	;

pattern_iterator: pattern_map_naked | pattern ;
\end{minted}

In its simplest form, the \code{for} statement is sued to iterate over lists,
maps or ranges.

\begin{minted}{\giovilexer}
for i in 0..10 do somethingWith(i)
for number in [1 2 3 4] do somethingWith(number)
for key in ["Simone"=>69 "Luca"=>420 "Paolo"=>10] do somethingWith(key)
\end{minted}

The for statement is like an usual for-each loop, but on steroids; the iterator
is actually a pattern (\see{ch:patterns}): for each iteration, the for
statement will match iterated value with the pattern, and perform the
appropriate deconstruction.

\begin{listing}
	\label{lst:for_loop_patterns_example}
	\caption{Example of various patterns used as iterators in \code{for} loops}
	\begin{minted}{\giovilexer}
for _ in 0..9 do something() # do something 10 times

for [[i]] in some2DMatrix do somethingWith(i)

mapOfNicknames [Text => [Text]] <- [
	"Simone" => ["Simon", "simonvic"]
	"Mario"  => ["Super Mario", "The plumber"]
	"Luigi"  => ["Green Mario"]
]
for [name => [nickname]] in mapOfNicknames do {
	somethingWith(name, nickname)
}

listOfColors [tuple<Int,Int,Int>] <- [
	<0 255 0>
	<55 0 255>
	<255 255 255>
]
for <r, g, b> in listOfColors do {
	"The color is made of $r $g $b" -> _out
}
	\end{minted}
\end{listing}

\nb{Nesting iterator patterns effectively means nesting iterations}

\todo{aggiungere RFC di iteration counters}


